\documentclass[11pt,a4paper]{article}

\usepackage[T1]{fontenc}
\usepackage[margin=2.5cm]{geometry}
\usepackage{amsmath,amssymb,amsthm}
\usepackage{bm}
\usepackage{graphicx}
\usepackage{hyperref}
\usepackage{enumitem}
\usepackage{booktabs}
\usepackage{xcolor}
\usepackage{tcolorbox}
\usepackage{parskip}
\usepackage{algorithm}
\usepackage{algorithmic}

% tcolorbox for examples
\tcbuselibrary{breakable}
\newtcolorbox{example}[1][]{
  colback=blue!3,
  colframe=blue!40!black,
  fonttitle=\bfseries,
  title=Numerical Example,
  breakable,
  #1
}

\newtcolorbox{keyidea}[1][]{
  colback=green!3,
  colframe=green!40!black,
  fonttitle=\bfseries,
  title=Key Idea,
  breakable,
  #1
}

\newtcolorbox{warning}[1][]{
  colback=red!3,
  colframe=red!40!black,
  fonttitle=\bfseries,
  title=Watch Out,
  breakable,
  #1
}

% Custom commands matching the paper
\newcommand{\R}{\mathbb{R}}
\newcommand{\N}{\mathcal{N}}
\newcommand{\E}{\mathbb{E}}
\newcommand{\G}{\mathcal{G}}
\newcommand{\bW}{\mathbf{W}}
\newcommand{\bx}{\mathbf{x}}
\newcommand{\bh}{\mathbf{h}}

\theoremstyle{definition}
\newtheorem{definition}{Definition}[section]
\theoremstyle{plain}
\newtheorem{theorem}{Theorem}[section]

\title{\textbf{Technical Companion to}\\[4pt]
\emph{Cost-Sensitive Neighborhood Aggregation\\for Heterophilous Graphs}\\[10pt]
{\large A Self-Contained Guide for CS Undergraduates}}

\author{Companion to the paper by Eyal Weiss}

\date{}

\begin{document}

\maketitle

\begin{abstract}
This document explains every equation and technical concept in the CSNA paper
in detail, assuming only undergraduate-level linear algebra, probability,
optimization, and graph theory.
No prior knowledge of graph neural networks, message passing, or GNN notation is
required.
Each equation is accompanied by (1) a plain-English description of what it computes,
(2) a symbol-by-symbol breakdown, (3) an intuitive explanation of why the
computation makes sense, and (4) a concrete numerical example where helpful.
\end{abstract}

\tableofcontents
\newpage

%% ================================================================
\section{Background: Graphs, Features, and the Classification Problem}
\label{sec:background}
%% ================================================================

\subsection{What Is a Graph in This Context?}

A \textbf{graph} $\G = (V, E)$ consists of:
\begin{itemize}
  \item A set of \textbf{nodes} $V$ with $n = |V|$ nodes, numbered $1, 2, \ldots, n$.
    Each node represents an entity---a webpage, a Wikipedia article, an actor, etc.
  \item A set of \textbf{edges} $E \subseteq V \times V$.
    An edge $(i, j) \in E$ means nodes $i$ and $j$ are connected (e.g., there is
    a hyperlink between two webpages).
    The graph is \emph{undirected}: if $(i,j) \in E$, then $(j,i) \in E$.
\end{itemize}

Every node $i$ carries two pieces of information:

\begin{enumerate}
  \item A \textbf{feature vector} $\bx_i \in \R^d$.
    This is a list of $d$ real numbers that describe the node.
    For webpages, the features might encode word frequencies; for actors, they
    might encode genre co-occurrences.
    We stack all feature vectors into a \textbf{feature matrix}
    $\mathbf{X} \in \R^{n \times d}$, where row $i$ is $\bx_i$.

  \item A \textbf{label} $y_i \in \{1, 2, \ldots, C\}$.
    This is the class or category of the node.
    The goal of \textbf{node classification} is to predict the label of each node
    using its features and the graph structure.
    During training, we know the labels of some nodes; at test time, we must
    predict the labels of the rest.
\end{enumerate}

\begin{example}
Suppose we have 4 nodes (webpages) with 3-dimensional features and 2 classes (``student page'' vs.\ ``faculty page''):

\begin{center}
\begin{tabular}{cccc}
\toprule
Node & Features $\bx_i$ & Label $y_i$ & Neighbors \\
\midrule
1 & $[1, 0, 3]$ & 1 (student) & $\{2, 3\}$ \\
2 & $[1, 0, 1]$ & 1 (student) & $\{1, 4\}$ \\
3 & $[0, 5, 2]$ & 2 (faculty) & $\{1, 4\}$ \\
4 & $[0, 4, 1]$ & 2 (faculty) & $\{2, 3\}$ \\
\bottomrule
\end{tabular}
\end{center}

Here $n = 4$, $d = 3$, $C = 2$.
Nodes 1 and 2 (students) are connected to each other,
but they are also connected to faculty nodes---this graph has significant heterophily.
\end{example}


\subsection{Edge Homophily Ratio $\mathcal{H}$}
\label{sec:homophily}

\begin{keyidea}
Homophily measures how often connected nodes share the same label.
High homophily means ``birds of a feather flock together''; low homophily
(heterophily) means connected nodes tend to differ.
\end{keyidea}

The \textbf{edge homophily ratio} is defined as:
\[
  \mathcal{H} = \frac{|\{(i,j) \in E : y_i = y_j\}|}{|E|}
\]

\textbf{What it computes:}
Count the number of edges whose endpoints share the same label, and divide by
the total number of edges.

\textbf{Symbol breakdown:}
\begin{itemize}[nosep]
  \item $\{(i,j) \in E : y_i = y_j\}$ --- the set of edges where both endpoints
    have the same label.
  \item $|\cdot|$ denotes set cardinality (the number of elements).
  \item $|E|$ is the total number of edges.
\end{itemize}

\textbf{Range and interpretation:}
\begin{itemize}[nosep]
  \item $\mathcal{H} = 1$: every edge connects same-label nodes (\textbf{perfect homophily}).
  \item $\mathcal{H} = 0$: every edge connects different-label nodes
    (\textbf{perfect heterophily}).
  \item $\mathcal{H} \approx 0.5$: edges are roughly random with respect to labels.
  \item $\mathcal{H} < 0.5$: the graph is \textbf{heterophilous}---connected
    nodes tend to differ in label.
\end{itemize}

\begin{example}
In the 4-node graph above, the edges are $\{(1,2), (1,3), (2,4), (3,4)\}$.
\begin{itemize}[nosep]
  \item Edge $(1,2)$: both students $\to$ same label.
  \item Edge $(1,3)$: student--faculty $\to$ different label.
  \item Edge $(2,4)$: student--faculty $\to$ different label.
  \item Edge $(3,4)$: both faculty $\to$ same label.
\end{itemize}
Same-label edges: 2 out of 4 total, so $\mathcal{H} = 2/4 = 0.5$.

Now suppose we remove edge $(1,2)$ and add edge $(1,4)$ instead.
Then all 4 edges are cross-class, giving $\mathcal{H} = 0/4 = 0$---perfect heterophily.

The datasets in the paper have $\mathcal{H}$ between 0.09 and 0.23, meaning that
fewer than 1 in 4 edges connect same-class nodes.
\end{example}


\subsection{Message Passing (Equation 1)}
\label{sec:message_passing}

\textbf{Message passing} is the fundamental operation of Graph Neural Networks (GNNs).
The idea is simple: to update node $i$'s representation, \emph{collect information
from its neighbors, combine it, and transform the result}.

The paper's Equation~1 defines one round of message passing:
\begin{equation}
  \bh_i^{(\ell+1)} = \sigma\!\Bigl(\sum_{j \in \N(i)} \alpha_{ij}\, \bW^{(\ell)}\, \bh_j^{(\ell)}\Bigr)
  \tag{Eq.~1}
\end{equation}

\textbf{What it computes in plain English:}
To get node $i$'s new representation at layer $\ell+1$, do the following: for each
neighbor $j$ (including $i$ itself), transform $j$'s current representation by the
weight matrix $\bW$, scale it by a weight $\alpha_{ij}$, sum all these contributions,
and pass the result through a nonlinear activation function.

\textbf{Symbol-by-symbol breakdown:}
\begin{itemize}
  \item $\bh_j^{(\ell)} \in \R^d$: the \textbf{representation} (also called
    ``embedding'' or ``hidden state'') of node $j$ at layer $\ell$.
    At layer 0, $\bh_j^{(0)} = \bx_j$ (the raw features).
    After each layer, representations are updated.

  \item $\bW^{(\ell)} \in \R^{d' \times d}$: a \textbf{learnable weight matrix}
    at layer $\ell$.
    This is the neural network's set of parameters---values that are adjusted
    during training to minimize classification error.
    It transforms a $d$-dimensional vector into a $d'$-dimensional vector.
    Multiplication $\bW^{(\ell)} \bh_j^{(\ell)}$ is a standard matrix--vector
    product, producing a vector in $\R^{d'}$.

  \item $\N(i)$: the \textbf{neighborhood} of node $i$, including $i$ itself
    (via a self-loop).
    If node $i$ has edges to nodes 2, 5, and 7, then $\N(i) = \{i, 2, 5, 7\}$.

  \item $\alpha_{ij} \in \R$: the \textbf{aggregation weight} for the message from
    $j$ to $i$.
    This controls how much influence neighbor $j$ has on node $i$'s updated
    representation.
    In the simplest GNN (GCN), $\alpha_{ij} = 1/\sqrt{d_i \cdot d_j}$ where $d_i$
    is the degree of node $i$---essentially an average.
    More sophisticated GNNs learn these weights (like attention in GAT).

  \item $\sum_{j \in \N(i)}$: the \textbf{aggregation step}---sum over all
    neighbors (including self).

  \item $\sigma(\cdot)$: a \textbf{nonlinear activation function}, such as
    ReLU ($\sigma(x) = \max(0, x)$, applied element-wise).
    Without nonlinearity, stacking multiple layers would collapse to a single
    linear transformation.
\end{itemize}

\begin{example}
Suppose node 1 has neighbors $\N(1) = \{1, 2, 3\}$ with representations:
\[
  \bh_1^{(0)} = [1, 0],\quad \bh_2^{(0)} = [0, 2],\quad \bh_3^{(0)} = [1, 1]
\]
and the weight matrix is $\bW = \begin{bmatrix} 1 & 0 \\ 0 & 1 \end{bmatrix}$
(identity, for simplicity),
with uniform weights $\alpha_{ij} = 1/3$ for all $j$.

Then:
\begin{align*}
  \bh_1^{(1)} &= \sigma\!\left(\tfrac{1}{3}\bW[1,0]^\top + \tfrac{1}{3}\bW[0,2]^\top
    + \tfrac{1}{3}\bW[1,1]^\top\right) \\
  &= \sigma\!\left(\tfrac{1}{3}[1,0] + \tfrac{1}{3}[0,2] + \tfrac{1}{3}[1,1]\right) \\
  &= \sigma\!\left([2/3,\; 1]\right) = [2/3,\; 1] \quad\text{(using ReLU)}
\end{align*}

The new representation of node 1 is the average of its neighbors' (transformed)
features, passed through ReLU.
\end{example}


\subsection{Why Homophily Makes Aggregation Work---and Heterophily Breaks It}
\label{sec:why_homo_works}

\begin{keyidea}
If your neighbors mostly share your label, averaging their features reinforces
your class identity.
If your neighbors mostly have a different label, averaging their features
\emph{dilutes} or \emph{overwrites} your class identity.
\end{keyidea}

Imagine you are a student node with features that ``look like a student.''
Your neighbors are:

\textbf{Homophilous case ($\mathcal{H}$ high):}
Most neighbors are also students.
Averaging their features produces something that still looks like a student---possibly
an even cleaner signal because noise averages out.
After aggregation, student nodes become \emph{more} similar to each other and
\emph{more} distinct from faculty nodes.
Classification is easy.

\textbf{Heterophilous case ($\mathcal{H}$ low):}
Most neighbors are faculty.
Averaging their features produces something that looks like a mix of student and
faculty---or even more like faculty than student.
After aggregation, your representation has been \emph{pulled toward} the other
class.
Student and faculty representations become harder to distinguish.
Classification suffers.

This is why, on highly heterophilous graphs (like Texas with $\mathcal{H} = 0.09$),
a simple Multi-Layer Perceptron (MLP) that \emph{ignores the graph entirely}
often outperforms standard GNNs: it is better to use only your own features
than to average in features from mostly-wrong-class neighbors.

The CSNA paper's central question is:
\emph{Can we design smarter aggregation weights that avoid this dilution?}


%% ================================================================
\section{The CSNA Method: Step by Step}
\label{sec:method}
%% ================================================================

CSNA (Cost-Sensitive Neighborhood Aggregation) modifies message passing by
assigning each edge an individualized weight based on how similar the two
endpoints look in a learned projection space.
We now walk through each step.


\subsection{Step 1: The Learned Projection $\bW_g$}
\label{sec:projection}

Before measuring how similar two nodes are, CSNA first \emph{projects} their
representations into a new space using a learned matrix $\bW_g$.

\textbf{Why not just compare the raw representations?}
Raw features live in a high-dimensional space (e.g., $d = 1703$ for the webpage
datasets) where many dimensions may be irrelevant to class identity.
Two students might differ wildly in raw features (one's page discusses algorithms,
another discusses databases), but in a well-chosen lower-dimensional projection,
the class-relevant signal is amplified and the irrelevant variation is discarded.

This is analogous to dimensionality reduction techniques like PCA, but here the
projection $\bW_g \in \R^{d' \times d}$ is \emph{learned end-to-end} via
backpropagation---the neural network discovers which projection makes same-class
nodes look similar and different-class nodes look different.


\subsection{Step 2: Edge Cost --- Observed Divergence $g_{ij}$ (Equation 2)}
\label{sec:gij}

\begin{equation}
  g_{ij}^{(\ell)} = \|\bW_g\, \bh_i^{(\ell)} - \bW_g\, \bh_j^{(\ell)}\|_2
  \tag{Eq.~2}
\end{equation}

\textbf{What it computes:}
The Euclidean (L2) distance between node $i$ and node $j$ after both are projected
into the learned space.

\textbf{Symbol breakdown:}
\begin{itemize}
  \item $\bh_i^{(\ell)} \in \R^d$: node $i$'s representation at layer $\ell$.
  \item $\bW_g \in \R^{d' \times d}$: the learned projection matrix.
    $d'$ is typically smaller than $d$ (e.g., 64 vs.\ 1703).
  \item $\bW_g \bh_i^{(\ell)} \in \R^{d'}$: node $i$'s projected representation.
  \item $\bW_g \bh_i^{(\ell)} - \bW_g \bh_j^{(\ell)} \in \R^{d'}$: the difference
    vector between the two projected representations.
  \item $\|\cdot\|_2$: the L2 (Euclidean) norm, i.e.,
    $\|\mathbf{v}\|_2 = \sqrt{v_1^2 + v_2^2 + \cdots + v_{d'}^2}$.
  \item $g_{ij}^{(\ell)} \geq 0$: a non-negative scalar. Low $g_{ij}$ means
    ``these nodes look similar in the projection'' (likely same class);
    high $g_{ij}$ means ``these nodes look different'' (likely different class).
\end{itemize}

\textbf{Why ``observed divergence''?}
The paper calls $g_{ij}$ the ``observed divergence'' because it is computed
directly from the observed (current) representations---no additional learned
scoring function is needed.
It is a direct measurement of how far apart $i$ and $j$ appear to be.

\begin{example}
Suppose $d' = 2$ and the projected representations are:
\[
  \bW_g \bh_1 = [3, 1], \quad \bW_g \bh_2 = [3, 0], \quad \bW_g \bh_3 = [0, 4]
\]
Then:
\begin{align*}
  g_{12} &= \|[3,1] - [3,0]\|_2 = \|[0, 1]\|_2 = 1 \quad\text{(small: similar)}\\
  g_{13} &= \|[3,1] - [0,4]\|_2 = \|[3, -3]\|_2 = \sqrt{9+9} = \sqrt{18} \approx 4.24 \quad\text{(large: different)}
\end{align*}

If nodes 1 and 2 are same-class and node 3 is a different class, the projection
has successfully made same-class nodes close and different-class nodes far apart.
\end{example}


\subsection{Step 3: The Sigmoid Concordance Score $s_{ij}$ (Equation 3)}
\label{sec:concordance}

\begin{equation}
  s_{ij}^{(\ell)} = \sigma\!\left(\frac{-g_{ij}^{(\ell)}}{\tau}\right) \in (0,\; 0.5]
  \tag{Eq.~3}
\end{equation}

\textbf{What it computes:}
Converts the distance $g_{ij}$ into a ``concordance score.''
Because $g_{ij} \geq 0$, the input to the sigmoid is always $-g_{ij}/\tau \leq 0$,
so $s_{ij} \leq \sigma(0) = 0.5$.
The maximum $s_{ij} = 0.5$ is attained only for self-loops ($g_{ii} = 0$);
for any edge connecting two distinct nodes with $g_{ij} > 0$, the score is strictly
less than $0.5$.
Higher scores (closer to $0.5$) indicate the edge is likely \textbf{concordant}
(same class); lower scores (closer to $0$) indicate the edge is likely
\textbf{discordant} (different class).

\textbf{Symbol breakdown:}
\begin{itemize}
  \item $g_{ij}^{(\ell)} \geq 0$: the distance from Eq.~2.
  \item $\tau > 0$: a \textbf{temperature} hyperparameter that controls the
    sharpness of the sigmoid.
    Tuned from $\{0.1, 0.5, 1.0, 2.0\}$ in experiments.
  \item $-g_{ij}/\tau$: a non-positive number. The minus sign means that
    \emph{larger distance $\to$ more negative input $\to$ lower sigmoid output}.
  \item $\sigma(z) = 1/(1 + e^{-z})$: the \textbf{sigmoid function},
    which maps any real number to $(0, 1)$.
\end{itemize}

\textbf{Why the sigmoid?}
The sigmoid provides a smooth, differentiable mapping from distance to a score,
which is needed for two reasons:
(1) we need a probabilistic-style weight to split messages between two channels,
and (2) the function must be differentiable so gradients can flow back to $\bW_g$
during training.
Note that because the sigmoid input $-g_{ij}/\tau$ is always non-positive,
the output is restricted to $(0, 0.5]$---not the full $(0,1)$ range of the sigmoid.

\textbf{What does temperature $\tau$ control?}
\begin{itemize}[nosep]
  \item Small $\tau$ (e.g., 0.1): the sigmoid transition is \emph{sharp}.
    Even a small distance produces $s_{ij} \approx 0$ (hard routing).
  \item Large $\tau$ (e.g., 2.0): the sigmoid transition is \emph{gradual}.
    Moderate distances still produce $s_{ij}$ near 0.5 (soft routing).
\end{itemize}
Think of $\tau$ as controlling how ``decisive'' the routing is.

\begin{example}
Continuing from the previous example with $\tau = 1$:
\begin{align*}
  s_{12} &= \sigma(-1/1) = \sigma(-1) = \frac{1}{1 + e^{1}} \approx \frac{1}{1 + 2.72} \approx 0.27 \\
  s_{13} &= \sigma(-4.24/1) = \sigma(-4.24) = \frac{1}{1 + e^{4.24}} \approx \frac{1}{1 + 69.4} \approx 0.014
\end{align*}

Node pair (1,2) gets concordance 0.27 (moderate---could be same class),
while pair (1,3) gets concordance 0.014 (very low---likely different class).

With $\tau = 0.1$: $s_{12} = \sigma(-10) \approx 0.00005$ and
$s_{13} = \sigma(-42.4) \approx 0$. Even the moderately-close pair gets near-zero
concordance---sharp temperature leaves little room for nuance.

With $\tau = 5$: $s_{12} = \sigma(-0.2) \approx 0.45$ and
$s_{13} = \sigma(-0.85) \approx 0.30$. Both scores are moderate---soft temperature
makes the routing very gentle.

Note: $s_{ii} = \sigma(0/\tau) = \sigma(0) = 0.5$ for self-loops, since $g_{ii} = 0$.
\end{example}


\subsection{Softmax Normalization of Weights}
\label{sec:softmax}

Before using the concordance scores as aggregation weights, they are normalized
using the \textbf{softmax} function, computed \emph{per node}:
\[
  \tilde{s}_{ij} = \mathrm{softmax}_{j \in \N(i)}(s_{ij})
  = \frac{\exp(s_{ij})}{\sum_{k \in \N(i)} \exp(s_{ik})}
\]
Similarly, the discordance weights use $1 - s_{ij}$:
\[
  \tilde{d}_{ij} = \mathrm{softmax}_{j \in \N(i)}(1 - s_{ij})
  = \frac{\exp(1 - s_{ij})}{\sum_{k \in \N(i)} \exp(1 - s_{ik})}
\]

\textbf{What softmax does here:}
For each node $i$, softmax converts the raw concordance scores of all its
neighbors into a probability distribution that sums to 1.
Neighbors with higher concordance get a larger share of the total weight in
the concordant channel; neighbors with higher discordance ($1 - s_{ij}$) get
a larger share in the discordant channel.

\textbf{Why normalize per node?}
Without normalization, a node with many neighbors would accumulate a much larger
aggregated signal than a node with few neighbors.
Per-node softmax ensures every node's aggregated representation is a weighted
\emph{average} (weights summing to 1) of its neighbors, regardless of degree.

\textbf{Why softmax instead of simple division?}
Softmax amplifies differences: if one neighbor has concordance 0.9 and
another has 0.2, softmax gives the first neighbor a disproportionately large
share (because $\exp(0.9) \gg \exp(0.2)$).
Simple normalization ($s_{ij} / \sum_k s_{ik}$) would produce more uniform weights.

\begin{example}
Suppose node 1 has three neighbors (including itself) with concordance scores:
\[
  s_{11} = 0.5, \quad s_{12} = 0.4, \quad s_{13} = 0.05
\]
Softmax for the concordant channel:
\begin{align*}
  \exp(0.5) &\approx 1.649, \quad \exp(0.4) \approx 1.492, \quad \exp(0.05) \approx 1.051 \\
  \text{Sum} &= 1.649 + 1.492 + 1.051 = 4.192 \\
  \tilde{s}_{11} &= 1.649/4.192 \approx 0.393 \\
  \tilde{s}_{12} &= 1.492/4.192 \approx 0.356 \\
  \tilde{s}_{13} &= 1.051/4.192 \approx 0.251
\end{align*}

Discordant channel uses $1 - s_{ij}$: values are $0.5, 0.6, 0.95$.
\begin{align*}
  \exp(0.5) &\approx 1.649, \quad \exp(0.6) \approx 1.822, \quad \exp(0.95) \approx 2.586 \\
  \text{Sum} &= 1.649 + 1.822 + 2.586 = 6.057 \\
  \tilde{d}_{11} &\approx 0.272, \quad
  \tilde{d}_{12} \approx 0.301, \quad
  \tilde{d}_{13} \approx 0.427
\end{align*}

Notice: in the concordant channel, the self-loop and same-class neighbor (node 2)
get the most weight.
In the discordant channel, the different-class neighbor (node 3) gets the most weight.
The channels emphasize different parts of the neighborhood.
\end{example}


\subsection{Step 4: Dual-Channel Aggregation (Equations 4--5)}
\label{sec:dual_channel}

\begin{align}
  \bh_i^{\text{con}} &= \sum_{j \in \N(i)} \tilde{s}_{ij}\, \bW_{\text{con}}\, \bh_j^{(\ell)}
  \tag{Eq.~4} \\[6pt]
  \bh_i^{\text{dis}} &= \sum_{j \in \N(i)} \tilde{d}_{ij}\, \bW_{\text{dis}}\, \bh_j^{(\ell)}
  \tag{Eq.~5}
\end{align}

\textbf{What they compute:}
Two separate aggregated representations for node $i$:
\begin{itemize}
  \item $\bh_i^{\text{con}}$ (\textbf{concordant channel}): a weighted sum of
    neighbors' features, where likely-same-class neighbors contribute more.
    Transformed by $\bW_{\text{con}}$.
  \item $\bh_i^{\text{dis}}$ (\textbf{discordant channel}): a weighted sum where
    likely-different-class neighbors contribute more.
    Transformed by a \emph{different} matrix $\bW_{\text{dis}}$.
\end{itemize}

\textbf{Symbol breakdown:}
\begin{itemize}
  \item $\tilde{s}_{ij}$: the softmax-normalized concordance weight (a scalar
    between 0 and 1, summing to 1 over $j$).
  \item $\tilde{d}_{ij}$: the softmax-normalized discordance weight.
  \item $\bW_{\text{con}} \in \R^{d' \times d}$: weight matrix for the concordant
    channel.
  \item $\bW_{\text{dis}} \in \R^{d' \times d}$: weight matrix for the discordant
    channel.
  \item $\bh_j^{(\ell)} \in \R^d$: neighbor $j$'s current representation.
  \item $\bh_i^{\text{con}}, \bh_i^{\text{dis}} \in \R^{d'}$: the resulting
    channel representations for node $i$.
\end{itemize}

\begin{keyidea}[title=Why Two DIFFERENT Weight Matrices Matter]
This is the heart of CSNA's design. The concordant and discordant channels use
\emph{independent} weight matrices $\bW_{\text{con}}$ and $\bW_{\text{dis}}$.

Why? Because same-class neighbors carry a fundamentally different type of
information than different-class neighbors:

\begin{itemize}[nosep]
  \item \textbf{Same-class neighbors} reinforce your identity. Their features,
    when transformed by $\bW_{\text{con}}$, should produce representations that
    sharpen your class signal. The concordant transformation might learn to
    extract and amplify class-specific features.

  \item \textbf{Different-class neighbors} carry complementary (or
    contradictory) information. Their features, when transformed by
    $\bW_{\text{dis}}$, might be used in a ``negative'' sense: ``if my
    different-class neighbors look like X, then I am probably not-X.''
    The discordant transformation can learn to extract useful contrast
    information.
\end{itemize}

If both channels used the same $\bW$, there would be no way to treat the two
types of information differently---the model would collapse to standard
weighted aggregation.
\end{keyidea}


\subsection{Step 5: The Gating Mechanism (Equation 6)}
\label{sec:gating}

\begin{equation}
  \bh_i^{(\ell+1)} = \sum_{k \in \{\text{con}, \text{dis}, \text{self}\}}
  \gamma_k(\bh_i) \cdot \bh_i^{k}
  \tag{Eq.~6}
\end{equation}
where:
\[
  \bh_i^{\text{self}} = \bW_{\text{self}}\, \bh_i^{(\ell)}
\]
\[
  \gamma(\bh_i) = \mathrm{softmax}\bigl(\bW_\gamma\, [\bh_i^{\text{con}} \| \bh_i^{\text{dis}} \| \bh_i^{\text{self}}]\bigr) \in \R^3
\]

\textbf{What it computes:}
The final updated representation of node $i$ is a weighted combination of three
ingredients:
\begin{enumerate}[nosep]
  \item $\bh_i^{\text{con}}$: what likely-same-class neighbors say (concordant channel).
  \item $\bh_i^{\text{dis}}$: what likely-different-class neighbors say (discordant channel).
  \item $\bh_i^{\text{self}}$: what the node itself says (ego channel, ignoring
    the graph).
\end{enumerate}
The gate weights $\gamma_{\text{con}}, \gamma_{\text{dis}}, \gamma_{\text{self}}$
determine the mixing proportions, and they \emph{depend on the node's own
representation}---different nodes can rely on different channels.

\textbf{Symbol breakdown:}
\begin{itemize}
  \item $\bW_{\text{self}} \in \R^{d' \times d}$: weight matrix for the ego
    (self) channel. This transforms the node's own representation without
    any neighbor information.
  \item $[\bh_i^{\text{con}} \| \bh_i^{\text{dis}} \| \bh_i^{\text{self}}] \in \R^{3d'}$:
    the concatenation of all three channel outputs.
    If each channel produces a 64-dimensional vector, this is a 192-dimensional
    vector.
  \item $\bW_\gamma \in \R^{3 \times 3d'}$: a small weight matrix that produces
    3 scores (one per channel) from the concatenated representation.
  \item $\mathrm{softmax}(\cdot) \in \R^3$: converts the 3 raw scores into a
    probability distribution: $\gamma_{\text{con}} + \gamma_{\text{dis}} +
    \gamma_{\text{self}} = 1$, with each $\gamma_k \geq 0$.
  \item The final output is: $\gamma_{\text{con}} \cdot \bh_i^{\text{con}} +
    \gamma_{\text{dis}} \cdot \bh_i^{\text{dis}} +
    \gamma_{\text{self}} \cdot \bh_i^{\text{self}}$.
\end{itemize}

\textbf{Why per-node gating?}
Different nodes benefit from different information sources.
A node surrounded by mostly-same-class neighbors should rely on the concordant
channel. A node with many different-class neighbors that carry useful contrast
information should lean toward the discordant channel. A node whose neighbors
are uninformative should fall back to its own features (self channel).
The per-node gate learns to make this decision for each node individually.

\textbf{Why softmax over 3 channels?}
Softmax ensures the gate weights are non-negative and sum to 1, so the output is
a convex combination of the three channel representations.
This prevents any channel from having negative influence (which could destabilize
training) and ensures the output scale stays bounded.

\begin{example}
Suppose for node $i$ the gate logits are $[1.0, -0.5, 2.0]$:
\begin{align*}
  \gamma &= \mathrm{softmax}([1.0, -0.5, 2.0]) \\
  &= \frac{1}{e^{1.0} + e^{-0.5} + e^{2.0}} [e^{1.0}, e^{-0.5}, e^{2.0}] \\
  &\approx \frac{1}{2.72 + 0.61 + 7.39}[2.72, 0.61, 7.39] \\
  &\approx [0.254, 0.057, 0.689]
\end{align*}

This node relies 69\% on its own features, 25\% on the concordant channel,
and only 6\% on the discordant channel---similar to an MLP with a small
graph-based correction.

On the Wisconsin and Actor datasets, the paper reports that the
\emph{average} gate weights across nodes are
$\gamma_{\text{self}} \approx 1.0$ (with $\gamma_{\text{con}} \approx 0$
and $\gamma_{\text{dis}} \approx 0$), indicating that the model
on average relies almost entirely on the ego channel, effectively
behaving like an MLP.
(These are averaged statistics, not per-node facts; individual nodes
may deviate from the average.)
\end{example}


\subsection{The Extended Variant: $f = g + h$ (Equation 7)}
\label{sec:extended}

\begin{equation}
  f_{ij}^{(\ell)} = g_{ij}^{(\ell)} + h_{ij}^{(\ell)}, \quad\text{where}\quad
  h_{ij}^{(\ell)} = \mathrm{softplus}\bigl(\mathbf{a}^\top [\bW_g \bh_i^{(\ell)} \| \bW_g \bh_j^{(\ell)}]\bigr)
  \tag{Eq.~7}
\end{equation}

\textbf{What it computes:}
An enriched edge cost that adds a \emph{learned asymmetric} component $h_{ij}$
on top of the observed divergence $g_{ij}$.
In the default ``lite'' variant, $f_{ij} = g_{ij}$ (no $h$ component).
In the ``extended'' variant, $f_{ij} = g_{ij} + h_{ij}$ replaces $g_{ij}$ in the
concordance formula (Eq.~3).

\textbf{Symbol breakdown:}
\begin{itemize}
  \item $g_{ij}$: the observed divergence from Eq.~2 (L2 distance in projection space).
  \item $[\bW_g \bh_i \| \bW_g \bh_j] \in \R^{2d'}$: the \textbf{concatenation}
    of $i$'s and $j$'s projected representations.
    If each projected representation is 64-dimensional, this is a 128-dimensional
    vector.
  \item $\mathbf{a} \in \R^{2d'}$: a learnable parameter vector.
  \item $\mathbf{a}^\top [\cdot] \in \R$: a dot product, producing a single scalar.
    This is essentially a learned linear scoring function on the pair $(i, j)$.
  \item $\mathrm{softplus}(z) = \ln(1 + e^z)$: a smooth, always-positive activation function.
    It ensures $h_{ij} \geq 0$, so that $h$ can only \emph{increase} the total
    cost (never decrease it below $g_{ij}$).
  \item $f_{ij} = g_{ij} + h_{ij} \geq g_{ij} \geq 0$: the combined cost.
\end{itemize}

\textbf{What does $h_{ij}$ add?}
The observed divergence $g_{ij}$ is symmetric ($g_{ij} = g_{ji}$) and purely
distance-based---it only knows \emph{how far apart} two nodes are.
The learned component $h_{ij}$ can capture more complex patterns by looking at
the actual values of the projected representations (not just their difference).
The concatenation $[\bW_g \bh_i \| \bW_g \bh_j]$ lets $\mathbf{a}$ learn
different roles for the source and target node features.

\textbf{Why softplus for non-negativity?}
The cost function should not be able to ``cancel out'' the observed divergence.
If $h_{ij}$ could be negative, it could make two distant nodes appear close
(cost near zero), defeating the purpose of measuring divergence.
Softplus ensures $h_{ij} > 0$, so the learned component can only add cost, not
subtract it.

\textbf{Practical impact:}
In the paper's experiments, the extended variant wins clearly only on Wisconsin
(84.1\% vs.\ 79.6\%) and is comparable or slightly worse on all other datasets.
The lite version (just $g_{ij}$) is the default because the added parameters
of $h_{ij}$ risk overfitting on the small datasets in the benchmark suite.


\subsection{Calibration Loss (Equation 8)}
\label{sec:calibration}

\begin{equation}
  \mathcal{L}_{\text{cal}} = \frac{1}{|E_{\text{train}}|}
  \sum_{(i,j) \in E_{\text{train}}}
  \bigl[\mathrm{ReLU}(g_{ij} - \mathbb{1}[y_i \neq y_j])\bigr]^2
  \tag{Eq.~8}
\end{equation}

\textbf{What it computes:}
A penalty term that encourages the cost function $g_{ij}$ to assign \emph{low}
cost to same-class edges.
It is added to the main classification loss during training.

\textbf{Symbol breakdown:}
\begin{itemize}
  \item $E_{\text{train}}$: the set of edges where both endpoints have known labels
    (i.e., both are in the training set).
  \item $\mathbb{1}[y_i \neq y_j]$: the \textbf{indicator function}---equals 1
    if $i$ and $j$ have different labels, and 0 if they have the same label.
  \item $g_{ij} - \mathbb{1}[y_i \neq y_j]$: the ``excess cost.''
    \begin{itemize}[nosep]
      \item For same-class edges ($y_i = y_j$): this equals $g_{ij} - 0 = g_{ij}$.
        Any positive cost is penalized.
      \item For different-class edges ($y_i \neq y_j$): this equals $g_{ij} - 1$.
        Only cost above 1 is penalized (cost up to 1 is ``free'').
    \end{itemize}
  \item $\mathrm{ReLU}(z) = \max(0, z)$: clips negative values to zero.
    Only \emph{excess} cost incurs a penalty.
  \item $[\cdot]^2$: squared penalty, so larger violations are penalized more heavily.
  \item The sum is averaged over the number of training edges.
\end{itemize}

\textbf{What it encourages:}
\begin{itemize}[nosep]
  \item Same-class edges: $g_{ij}$ should be as close to 0 as possible
    (these nodes should look similar in the projection).
  \item Different-class edges: $g_{ij}$ can be up to 1 without penalty,
    and is only penalized if it exceeds 1 (preventing extreme costs that
    could cause numerical issues).
\end{itemize}

\begin{warning}[title=The Scale Mismatch]
The indicator $\mathbb{1}[y_i \neq y_j]$ is either 0 or 1, while $g_{ij}$
(an L2 distance in a learned space) can take any non-negative value---potentially
much larger than 1.
This means the regularizer acts as a ``soft'' penalty rather than a precise
calibration: it pushes same-class costs toward zero but does not enforce
a specific threshold for different-class edges.
The paper explicitly notes this scale mismatch.
\end{warning}

The full training objective is:
\[
  \mathcal{L} = \mathcal{L}_{\text{CE}} + \lambda_{\text{cal}} \cdot \mathcal{L}_{\text{cal}}
\]
where $\mathcal{L}_{\text{CE}}$ is the standard cross-entropy classification loss and
$\lambda_{\text{cal}} = 0.1$ is fixed across all experiments.

\begin{example}
Suppose we have three training edges with costs and labels:
\begin{center}
\begin{tabular}{cccc}
\toprule
Edge & $g_{ij}$ & Same class? & Penalty contribution \\
\midrule
$(1,2)$ & 0.3 & Yes ($\mathbb{1}=0$) & $\mathrm{ReLU}(0.3 - 0)^2 = 0.09$ \\
$(1,3)$ & 2.5 & No ($\mathbb{1}=1$) & $\mathrm{ReLU}(2.5 - 1)^2 = 2.25$ \\
$(2,3)$ & 0.8 & No ($\mathbb{1}=1$) & $\mathrm{ReLU}(0.8 - 1)^2 = \mathrm{ReLU}(-0.2)^2 = 0$ \\
\bottomrule
\end{tabular}
\end{center}

$\mathcal{L}_{\text{cal}} = (0.09 + 2.25 + 0) / 3 = 0.78$.

The penalty comes from two sources: edge $(1,2)$ is same-class but has cost 0.3
(should be near 0), and edge $(1,3)$ is different-class but has cost 2.5 (above
the ``free'' threshold of 1).
Edge $(2,3)$ is different-class with cost 0.8, which is below 1, so no penalty.
\end{example}


\subsection{Architecture Choices}
\label{sec:architecture}

The paper makes several practical design choices that are worth understanding:

\paragraph{Input MLP.}
Before the first CSNA layer, the raw features $\mathbf{X}$ are passed through a
standard Multi-Layer Perceptron (2-layer feedforward network).
This projects the possibly very high-dimensional features ($d = 1703$) into a
lower-dimensional space, producing the initial representations $\bh_i^{(0)}$.
The MLP also introduces nonlinearity, which can help separate classes even before
any graph information is used.

\paragraph{Residual connections.}
Between CSNA layers, a \textbf{residual connection} adds the input to the output:
$\bh_i^{(\ell+1)} = \text{CSNA}(\bh_i^{(\ell)}) + \bh_i^{(\ell)}$.
This is the same idea as in ResNets: it makes it easy for the network to learn
the identity function (``do nothing'') if an additional layer would hurt.
It also helps with gradient flow during training.

\paragraph{Gate bias initialization $[0, 0, 1]$.}
The gate bias for $\bW_\gamma$ is initialized to $[0, 0, 1]$, which means:
\[
  \gamma = \mathrm{softmax}([0, 0, 1]) \approx [0.21, 0.21, 0.58]
\]
This favors the \textbf{self (ego) channel} at initialization.
The model starts behaving approximately like an MLP, which is a reasonable
baseline on heterophilous graphs.
During training, the gate can learn to shift weight toward the concordant or
discordant channels if the graph structure is useful.

Starting from MLP-like behavior is a deliberate safety measure: if the graph
structure turns out to be harmful (as on some heterophilous datasets), the model
does not need to ``unlearn'' reliance on bad graph information---it was never
relying on it in the first place.


\subsection{Algorithm 1: The Full CSNA Layer}
\label{sec:algorithm}

The paper consolidates all five steps described above into a single algorithm.
We reproduce it here for reference.
This is the \emph{lite} version (using only the observed divergence $g_{ij}$;
for the extended variant, replace $g_{ij}$ with $f_{ij} = g_{ij} + h_{ij}$).

\begin{algorithm}[h]
\caption{CSNA Layer (Lite Version)}
\begin{algorithmic}[1]
\REQUIRE Node features $\mathbf{H}^{(\ell)} \in \R^{n \times d}$, edge index $E$, temperature $\tau$
\ENSURE Updated features $\mathbf{H}^{(\ell+1)} \in \R^{n \times d'}$
\FOR{each edge $(i,j) \in E$}
  \STATE $g_{ij} \leftarrow \|\bW_g \bh_i - \bW_g \bh_j\|_2$ \hfill \{Pairwise distance in learned projection\}
  \STATE $s_{ij} \leftarrow \sigma(-g_{ij}/\tau)$ \hfill \{Concordance score\}
\ENDFOR
\FOR{each node $i \in V$}
  \STATE $\bh_i^{\text{con}} \leftarrow \sum_{j \in \N(i)} \tilde{s}_{ij}\, \bW_{\text{con}}\, \bh_j$ \hfill \{Concordant channel\}
  \STATE $\bh_i^{\text{dis}} \leftarrow \sum_{j \in \N(i)} \tilde{d}_{ij}\, \bW_{\text{dis}}\, \bh_j$ \hfill \{Discordant channel\}
  \STATE $\bh_i^{\text{self}} \leftarrow \bW_{\text{self}}\, \bh_i$ \hfill \{Ego transform\}
  \STATE $\bh_i^{(\ell+1)} \leftarrow \sum_k \gamma_k(\bh_i) \cdot \bh_i^k$ \hfill \{Gated combination\}
\ENDFOR
\end{algorithmic}
\end{algorithm}

\textbf{Reading the algorithm:}
\begin{itemize}
  \item \textbf{Lines 1--4} (edge loop): For every edge in the graph, compute the
    observed divergence $g_{ij}$ (Eq.~2) and convert it to a concordance score
    $s_{ij}$ (Eq.~3). Self-loops are included; since $g_{ii} = 0$, self-loops
    receive $s_{ii} = 0.5$ (the maximum concordance).
  \item \textbf{Lines 5--10} (node loop): For each node, compute the three channel
    representations using softmax-normalized weights
    ($\tilde{s}_{ij}$ and $\tilde{d}_{ij}$, from Section~\ref{sec:softmax}),
    then combine them via the learned per-node gate $\gamma$ (Eq.~6).
  \item The algorithm is applied $L$ times (the paper uses $L=2$), with residual
    connections between layers.
\end{itemize}


%% ================================================================
\section{Theoretical Analysis: When Does Cost-Sensitive Aggregation Help?}
\label{sec:theory}
%% ================================================================

The theoretical section of the paper asks: \emph{under what conditions does
weighting edges by their cost function improve classification over simple
mean aggregation?}
To answer this precisely, the analysis uses a synthetic graph model.


\subsection{The Contextual Stochastic Block Model (CSBM)}
\label{sec:csbm}

\begin{definition}[Contextual Stochastic Block Model]
A graph $\G$ is drawn from $\mathrm{CSBM}(n, 2, p, q, \mu)$:
\begin{itemize}[nosep]
  \item $n$ nodes split into two equal classes: $V_+$ (class +1) and $V_-$ (class $-1$),
    each of size $n/2$.
  \item Edges are drawn independently at random:
    \begin{itemize}[nosep]
      \item Probability $p$ for an edge between two same-class nodes.
      \item Probability $q$ for an edge between two different-class nodes.
    \end{itemize}
  \item Node features: $\bx_i \sim \mathcal{N}(\boldsymbol{\mu}_{y_i}, \mathbf{I}_d)$.
    \begin{itemize}[nosep]
      \item Class $+1$ nodes: features are drawn from a Gaussian centered at
        $\boldsymbol{\mu}_+ = +\frac{\mu}{2}\mathbf{e}_1$.
      \item Class $-1$ nodes: features are drawn from a Gaussian centered at
        $\boldsymbol{\mu}_- = -\frac{\mu}{2}\mathbf{e}_1$.
    \end{itemize}
\end{itemize}
\end{definition}

\textbf{What does this model capture?}
\begin{itemize}
  \item The \textbf{graph structure} is controlled by $p$ and $q$:
    \begin{itemize}[nosep]
      \item $p > q$: homophilous graph (same-class nodes more likely to connect).
      \item $p < q$: heterophilous graph (different-class nodes more likely to connect).
      \item $p = q$: random graph (edges independent of class).
    \end{itemize}
  \item The \textbf{node features} are noisy observations of class identity.
    Each node's feature vector is a random draw from a Gaussian whose center
    depends on the class. The parameter $\mu$ controls the \textbf{signal
    strength}: large $\mu$ means the two classes have very different mean features;
    small $\mu$ means the classes are hard to distinguish from features alone.
  \item $\mathbf{e}_1 = [1, 0, 0, \ldots, 0]$ is the first standard basis vector.
    The class means differ only in the first coordinate: $\boldsymbol{\mu}_+ -
    \boldsymbol{\mu}_- = \mu \cdot \mathbf{e}_1$.
  \item The homophily ratio in this model is $\mathcal{H} = p/(p+q)$.
\end{itemize}

\begin{example}
With $n = 100$, $p = 0.1$, $q = 0.4$, $\mu = 2$:
\begin{itemize}[nosep]
  \item 50 nodes per class.
  \item Each node connects to $\sim 0.1 \times 49 \approx 5$ same-class nodes
    and $\sim 0.4 \times 50 = 20$ different-class nodes.
  \item Homophily: $\mathcal{H} = 0.1/0.5 = 0.2$---highly heterophilous.
  \item Features: class +1 nodes have features centered at $[+1, 0, \ldots, 0]$;
    class $-1$ nodes at $[-1, 0, \ldots, 0]$.
  \item A typical class +1 node might have features $[1.3, -0.5, 0.2, \ldots]$
    (noisy version of $[1, 0, 0, \ldots]$).
\end{itemize}
\end{example}


\subsection{Theorem 1: Signal Distortion Under Mean Aggregation}
\label{sec:theorem1}

\begin{theorem}[Signal distortion under mean aggregation]
Let $\G \sim \mathrm{CSBM}(n, 2, p, q, \mu)$ with constant $p, q \in (0,1)$ and $q > p$ (heterophilous).
After one round of symmetrically normalized aggregation with self-loops, the expected class-mean difference satisfies:
\begin{equation}
  \E[\bar{\bh}_+^{(1)} - \bar{\bh}_-^{(1)}]
  = \frac{p - q}{p + q}(\boldsymbol{\mu}_+ - \boldsymbol{\mu}_-) + O(1/n)
  \tag{Eq.~9, Thm.~1}
\end{equation}
\end{theorem}

\textbf{What it says in plain English:}
After mean aggregation, the class-mean difference is scaled by the signed factor $\lambda = (p-q)/(p+q)$.
When $q > p$, this factor is \emph{negative}, meaning the class centroids swap sides---their relative direction \emph{reverses}.
The magnitude $|\lambda|$ is less than 1 (attenuation) whenever $q \neq p$, and is minimized near $p = q$---not in the extreme heterophily limit.
In fact, when $q \gg p$, $|\lambda| \to 1$: strong heterophily distorts the signal (sign flip) but does not collapse it.

\textbf{Symbol breakdown:}
\begin{itemize}
  \item $\bar{\bh}_+^{(1)}$: the mean representation of all class +1 nodes
    after one layer of aggregation.
    $\bar{\bh}_+^{(1)} = \frac{2}{n}\sum_{i \in V_+} \bh_i^{(1)}$.
  \item $\bar{\bh}_-^{(1)}$: same for class $-1$.
  \item $\bar{\bh}_+^{(1)} - \bar{\bh}_-^{(1)}$: the vector difference
    between the two class centroids after aggregation.
  \item $\boldsymbol{\mu}_+ - \boldsymbol{\mu}_-$: the original class-mean
    difference \emph{before} aggregation.
  \item $\frac{p-q}{p+q}$: the \textbf{scaling factor}. It is \emph{signed}:
    negative when $q > p$ (heterophily), positive when $p > q$ (homophily).
    When negative, the class centroids reverse direction.
  \item $O(1/n)$: a correction term that absorbs the self-loop contribution
    (negligible under the dense CSBM where degrees are $\Theta(n)$) and
    concentration errors.
\end{itemize}

\textbf{Intuitive explanation:}
Under mean aggregation, each node's new representation is an average of its
neighbors' features.
In a CSBM with $q > p$, a typical node has:
\begin{itemize}[nosep]
  \item Fraction $\frac{p}{p+q}$ of same-class neighbors (contributing
    toward $\boldsymbol{\mu}_+$).
  \item Fraction $\frac{q}{p+q}$ of different-class neighbors (contributing
    toward $\boldsymbol{\mu}_-$).
\end{itemize}

The average representation of a class +1 node becomes:
\[
  \frac{p}{p+q}\boldsymbol{\mu}_+ + \frac{q}{p+q}\boldsymbol{\mu}_-
  = \frac{p-q}{p+q}\boldsymbol{\mu}_+
\]
(using $\boldsymbol{\mu}_- = -\boldsymbol{\mu}_+$ in the signal direction).

The class signal $\boldsymbol{\mu}_+$ has been multiplied by
$\frac{p-q}{p+q}$, which is negative (since $q > p$) and has absolute value less
than 1.
The signal has both \emph{shrunk} and \emph{flipped sign}!

\begin{example}
With $p = 0.1$, $q = 0.4$:
\[
  \text{Attenuation factor} = \left(\frac{0.1 - 0.4}{0.1 + 0.4}\right)^2
  = \left(\frac{-0.3}{0.5}\right)^2 = (-0.6)^2 = 0.36
\]

The between-class separation shrinks to 36\% of the original.
Starting from $\mu^2 = 4$ (using $\mu = 2$), after aggregation:
$\E[\|\bar{\bh}_+ - \bar{\bh}_-\|^2] \approx 0.36 \times 4 = 1.44$.

With extreme heterophily ($p = 0.01$, $q = 0.99$):
\[
  \left(\frac{0.01 - 0.99}{0.01 + 0.99}\right)^2 = (-0.98)^2 = 0.9604
\]

Wait---this is close to 1! The signal is \emph{almost preserved}?
Yes, but the sign has flipped: class +1 nodes now have representations that
look like class $-1$ and vice versa. The $(\cdot)^2$ hides the sign flip.
The linear factor is $\frac{p-q}{p+q} = -0.98$, meaning the class centroids
have essentially swapped positions.
A classifier can still separate them, but a naive multi-layer GNN would compound
the problem at each layer.

The worst case is $q = p$ (random graph): the attenuation factor is 0,
meaning the class signal is completely destroyed.
\end{example}


\subsection{Theorem 2: Signal Preservation Under Cost-Sensitive Aggregation}
\label{sec:theorem2}

\begin{theorem}[Cost-sensitive aggregation preserves signal direction]
Under the same CSBM, suppose an aggregation scheme weights edge $(i,j)$ by $w_{ij}$,
with:
\begin{itemize}[nosep]
  \item $\E[w_{ij} \mid y_i = y_j] = w_+$ (average weight on same-class edges).
  \item $\E[w_{ij} \mid y_i \neq y_j] = w_-$ (average weight on different-class edges).
  \item $w_+ > w_- \geq 0$ (same-class edges get higher average weight).
\end{itemize}
Then:
\begin{equation}
  \E[\bar{\bh}_+^{(1)} - \bar{\bh}_-^{(1)}]
  = \frac{p\, w_+ - q\, w_-}{p\, w_+ + q\, w_-}(\boldsymbol{\mu}_+ - \boldsymbol{\mu}_-) + O(1/n)
  \tag{Eq.~11, Thm.~2}
\end{equation}

The scaling factor is positive---preserving the original class direction---if and only if $w_+/w_- > q/p$.
\end{theorem}

\textbf{What it says in plain English:}
If you can assign higher weights to same-class edges and lower weights to
different-class edges, you can control whether the class-mean direction is
preserved or reversed after aggregation.
The critical condition is that your weight ratio ($w_+/w_-$) must exceed the
graph's heterophily ratio ($q/p$) to keep the class centroids on the correct side.

\textbf{Symbol breakdown:}
\begin{itemize}
  \item $w_+$: the expected weight assigned to same-class edges. In CSNA,
    this corresponds to the average concordance score $s_{ij}$ for same-class pairs.
  \item $w_-$: the expected weight assigned to different-class edges.
  \item $p\, w_+$: the ``effective contribution'' of same-class edges
    (edge probability $\times$ weight).
  \item $q\, w_-$: the ``effective contribution'' of different-class edges.
  \item The formula replaces $p \to p\, w_+$ and $q \to q\, w_-$ in Theorem~1's
    formula. Cost-sensitive weighting effectively ``corrects'' the edge probabilities
    by reweighting them.
\end{itemize}

\textbf{The condition $w_+/w_- > q/p$---why it is intuitive:}

\begin{keyidea}[title=The Key Condition: $w_+/w_- > q/p$]
The graph's heterophily is captured by $q/p$: this is how much more likely a
different-class edge is compared to a same-class edge.
For example, if $q/p = 4$, there are 4 times as many cross-class edges as
same-class edges.

Your weighting scheme's discrimination is captured by $w_+/w_-$: how much more
weight you assign to same-class edges compared to different-class edges.

The condition says: \textbf{your weighting must discriminate edges at least as strongly as
the graph's heterophily confuses them}, in order to preserve the correct class direction.

If $q/p = 4$, you need $w_+/w_- > 4$---you must assign same-class edges more
than 4 times the weight of different-class edges.
The cross-class edges outnumber same-class edges 4-to-1, so your weights must
compensate by favoring same-class edges by more than 4-to-1.

When the condition holds, the scaling factor is positive: the class centroids stay
on the correct side.
When it fails ($w_+/w_- \leq q/p$), the signal direction reverses, just as
in unweighted mean aggregation.
\end{keyidea}

\begin{example}
With $p = 0.1$, $q = 0.4$ (so $q/p = 4$):

\textbf{Mean aggregation} ($w_+ = w_- = 1$):
\[
  \text{Factor} = \left(\frac{0.1 \cdot 1 - 0.4 \cdot 1}{0.1 \cdot 1 + 0.4 \cdot 1}\right)^2
  = \left(\frac{-0.3}{0.5}\right)^2 = 0.36
\]

\textbf{Cost-sensitive with $w_+ = 5, w_- = 1$} (ratio $w_+/w_- = 5 > 4 = q/p$):
\[
  \text{Factor} = \left(\frac{0.1 \cdot 5 - 0.4 \cdot 1}{0.1 \cdot 5 + 0.4 \cdot 1}\right)^2
  = \left(\frac{0.5 - 0.4}{0.5 + 0.4}\right)^2
  = \left(\frac{0.1}{0.9}\right)^2 \approx 0.012
\]

Wait---0.012 is \emph{less} than 0.36! How is this an improvement?

The key is the sign: with mean aggregation, $\frac{p-q}{p+q} = -0.6$
(the signal \emph{flips}). With cost-sensitive weighting,
$\frac{pw_+ - qw_-}{pw_+ + qw_-} = +0.11$ (the signal stays in the
\emph{correct direction}, though it is small).

For classification, the direction matters: a flipped signal at $-0.6$ is
actually usable (a classifier can learn the flip), but the multi-layer problem
is that repeated sign flips across layers cause oscillation.
Cost-sensitive weighting avoids this by keeping the signal in the correct direction.

\textbf{Better cost-sensitive weighting with $w_+ = 10, w_- = 1$}:
\[
  \text{Factor} = \left(\frac{1.0 - 0.4}{1.0 + 0.4}\right)^2
  = \left(\frac{0.6}{1.4}\right)^2 \approx 0.184
\]

Now the signal retains 18.4\% of its original strength in the correct direction,
compared to 36\% in the wrong direction under mean aggregation.
Stronger discrimination ($w_+/w_- = 10$) recovers more of the signal.

\textbf{Perfect oracle with $w_+ = 1, w_- = 0$} (ignore all cross-class edges):
\[
  \text{Factor} = \left(\frac{0.1 \cdot 1 - 0.4 \cdot 0}{0.1 \cdot 1 + 0.4 \cdot 0}\right)^2
  = \left(\frac{0.1}{0.1}\right)^2 = 1.0
\]

Perfect preservation! If you could perfectly identify and ignore all cross-class
edges, aggregation would not degrade the signal at all.
\end{example}


\subsection{Connecting the Theory to CSNA}
\label{sec:theory_connection}

\textbf{How does this apply to CSNA specifically?}

Theorem~2 applies to CSNA's concordant channel with:
\begin{itemize}[nosep]
  \item $w_{ij} = s_{ij}$ (the concordance score).
  \item $w_+ = \E[s_{ij} \mid y_i = y_j]$ (average concordance on same-class edges).
  \item $w_- = \E[s_{ij} \mid y_i \neq y_j]$ (average concordance on different-class edges).
\end{itemize}

For the theorem's sign-preservation condition to hold, CSNA's learned cost function
must assign concordance scores such that $w_+/w_- > q/p$.
The \textbf{calibration regularizer} (Eq.~8) encourages this by penalizing high costs
on same-class edges, which pushes $s_{ij}$ higher for same-class pairs and thus
increases $w_+/w_-$.

\textbf{Important caveat:}
Theorem~2 is \emph{conditional}---it says what happens \emph{if} the cost function
achieves good separation ($w_+/w_- > q/p$), but it does not guarantee that CSNA
\emph{will} learn such separation.
Also note that the theorem bounds the class-mean difference, not classification accuracy:
a sign-reversed but well-separated representation can still be corrected by a
downstream linear classifier.
Whether CSNA can actually learn a good cost function depends on the data: on
adversarial-heterophily datasets (Texas, Wisconsin, Cornell), the paper shows
that CSNA does learn good separation (high AUC for edge-type classification);
on informative-heterophily datasets (Chameleon, Squirrel), it does not.


\subsection{Proof Details: Equations from the Appendix}
\label{sec:proof_details}

The paper's appendix provides detailed proofs of Theorems~1 and~2.
We walk through the key intermediate equations here.

\subsubsection{Proof of Theorem 1 (Appendix E.1)}

Consider the binary CSBM with $n$ nodes, two equal classes $V_+$ and $V_-$,
intra-class edge probability $p$, and inter-class edge probability $q > p$.
For a node $i \in V_+$, the number of same-class neighbors is
$n_s(i) \sim \mathrm{Bin}(n/2 - 1, p)$ and different-class neighbors
$n_d(i) \sim \mathrm{Bin}(n/2, q)$.
Under the dense CSBM (degrees are $\Theta(n)$), the self-loop in
$\mathbf{A} + \mathbf{I}$ contributes $O(1/n)$ per node and does not affect
the leading coefficient.

The expected post-aggregation representation for a class $+1$ node is
(paper Equation~10, also labeled~13 in the appendix proof):
\begin{equation}
  \E[\bh_i^{(1)} | y_i = +1]
  = \frac{p}{p+q}\,\boldsymbol{\mu}_+
  + \frac{q}{p+q}\,\boldsymbol{\mu}_-
  + O(1/n)
  = \frac{p - q}{p + q} \cdot \frac{\mu}{2}\,\mathbf{e}_1 + O(1/n).
  \tag{Eq.~10}
\end{equation}

\textbf{What it says:}
Each class $+1$ node's expected representation after one round of mean aggregation
is a weighted average of the two class means, where the weights are the
relative edge probabilities.
The fraction $p/(p+q)$ of its neighbors are same-class (contributing
$\boldsymbol{\mu}_+$) and $q/(p+q)$ are different-class (contributing
$\boldsymbol{\mu}_-$).
Since $\boldsymbol{\mu}_+ = +\frac{\mu}{2}\mathbf{e}_1$ and
$\boldsymbol{\mu}_- = -\frac{\mu}{2}\mathbf{e}_1$, the combination
simplifies to $(p-q)/(p+q) \cdot \frac{\mu}{2}\,\mathbf{e}_1$.
The $O(1/n)$ term absorbs the self-loop contribution and the concentration
error (by Hoeffding's inequality, $n_s(i)/d_i$ concentrates around $p/(p+q)$
at rate $O(1/\sqrt{n})$, giving $O(1/n)$ after averaging over $n/2$ nodes per class).

By symmetry, $\E[\bh_i^{(1)} | y_i = -1] = \frac{p-q}{p+q} \cdot (-\frac{\mu}{2})\,\mathbf{e}_1 + O(1/n)$.
Averaging over classes yields the class-mean separation (Equation~14 in the appendix):
\begin{equation}
  \E[\bar{\bh}_+^{(1)} - \bar{\bh}_-^{(1)}]
  = \frac{p - q}{p + q}(\boldsymbol{\mu}_+ - \boldsymbol{\mu}_-) + O(1/n).\;\square
  \tag{Eq.~14}
\end{equation}

This completes the proof of Theorem~1: the scaling factor $\lambda = (p-q)/(p+q)$
is negative when $q > p$, confirming that mean aggregation reverses the
label-aligned signal direction under heterophily.


\subsubsection{Proof of Theorem 2 (Appendix E.2)}

Now suppose the aggregation assigns expected weight $w_+$ to same-class edges
and $w_-$ to different-class edges.
The self-loop still contributes $O(1/n)$ as in Theorem~1.
At leading order, the expected representation of a class $+1$ node becomes
(paper Equation~12, also labeled~15 in the appendix):
\begin{equation}
  \E[\bh_i^{(1)} | y_i = +1]
  = \frac{p\, w_+}{p\, w_+ + q\, w_-}\,\boldsymbol{\mu}_+
  + \frac{q\, w_-}{p\, w_+ + q\, w_-}\,\boldsymbol{\mu}_-
  + O(1/n).
  \tag{Eq.~12}
\end{equation}

\textbf{What changed from Theorem~1:}
Compared with Eq.~10, the edge probabilities $p$ and $q$ are now replaced by
\emph{effective} contributions $p\,w_+$ and $q\,w_-$.
If the weighting scheme assigns high weight to same-class edges ($w_+ \gg w_-$),
the effective contribution of same-class neighbors is amplified, counteracting
the numerical dominance of cross-class edges.

Setting $w_+ = w_- = 1$ recovers Eq.~10 (Theorem~1).

Averaging over classes (paper Equation~16 in the appendix):
\begin{equation}
  \E[\bar{\bh}_+^{(1)} - \bar{\bh}_-^{(1)}]
  = \frac{p\, w_+ - q\, w_-}{p\, w_+ + q\, w_-}
  (\boldsymbol{\mu}_+ - \boldsymbol{\mu}_-) + O(1/n).
  \tag{Eq.~16}
\end{equation}

The numerator $p\, w_+ - q\, w_-$ is positive if and only if
$w_+/w_- > q/p$. This completes the proof. $\square$

\begin{example}
Why does the condition $w_+/w_- > q/p$ arise?

The number of same-class edges scales with $p$, and the number of cross-class
edges scales with $q$.
In the aggregated sum, same-class edges contribute a total ``mass'' proportional
to $p \cdot w_+$ pulling toward $\boldsymbol{\mu}_+$, while cross-class edges
contribute mass proportional to $q \cdot w_-$ pulling toward
$\boldsymbol{\mu}_-$.
For the positive class direction to win, we need
$p \cdot w_+ > q \cdot w_-$, i.e., $w_+/w_- > q/p$.

With $p = 0.1$ and $q = 0.4$: the condition requires $w_+/w_- > 4$.
If CSNA learns $w_+ = 0.45$ and $w_- = 0.10$ (concordance scores for
same-class and different-class edges), then $w_+/w_- = 4.5 > 4$---the
condition is satisfied and the signal direction is preserved.
\end{example}


\subsubsection{Appendix Equations 13, 15, and 17}

For completeness, we list the appendix-specific equation labels that
correspond to the derivations above.
The appendix uses its own numbering for detailed proof steps:

\textbf{Equation 13} (Appendix E.1) gives the per-node expected representation
under mean aggregation, identical in content to Eq.~10 above:
\begin{equation}
  \E[\bh_i^{(1)} | y_i = +1]
  = \frac{p}{p+q}\,\boldsymbol{\mu}_+
  + \frac{q}{p+q}\,\boldsymbol{\mu}_-
  + O(1/n)
  = \frac{p - q}{p + q} \cdot \frac{\mu}{2}\,\mathbf{e}_1 + O(1/n).
  \tag{Eq.~13}
\end{equation}

\textbf{Equation 15} (Appendix E.2) gives the per-node expected representation
under cost-sensitive aggregation, identical in content to Eq.~12 above:
\begin{equation}
  \E[\bh_i^{(1)} | y_i = +1]
  = \frac{p\, w_+}{p\, w_+ + q\, w_-}\,\boldsymbol{\mu}_+
  + \frac{q\, w_-}{p\, w_+ + q\, w_-}\,\boldsymbol{\mu}_-
  + O(1/n).
  \tag{Eq.~15}
\end{equation}

\textbf{Equation 17} (Appendix G) decomposes the individual node's deviation
from the class mean into feature noise and edge-composition terms, and is
discussed in Section~\ref{sec:within_class_scatter} below.


\subsection{Within-Class Scatter (Appendix G)}
\label{sec:within_class_scatter}

Theorem~2 bounds the \emph{between-class} separation
$\E[\bar{\bh}_+^{(1)} - \bar{\bh}_-^{(1)}]$, but classification quality also
depends on how tightly clustered nodes are \emph{within} each class.
Appendix~G of the paper addresses this by analyzing the within-class scatter
matrix $\mathbf{S}_W$.

\textbf{The decomposition.}
For a node $i \in V_+$, the deviation from the class mean can be written as
(paper Equation~17):
\begin{equation}
  \bh_i^{(1)} - \bar{\bh}_+^{(1)}
  = \sum_{j \in \N(i)} w_{ij}\bigl(\bx_j - \E[\bx_j | y_j]\bigr)
  + \text{(edge-composition terms)}.
  \tag{Eq.~17}
\end{equation}

The first term captures \textbf{feature noise}: each neighbor $j$'s feature
vector deviates from the class-conditional mean $\E[\bx_j | y_j]$ due to the
random draw $\bx_j \sim \mathcal{N}(\boldsymbol{\mu}_{y_j}, \mathbf{I}_d)$.
This noise is averaged over $d_i = \Theta(n)$ neighbors, giving variance
$O(d/n)$.

The edge-composition terms capture the fact that different nodes have
different numbers of same-class and different-class neighbors
(random variation around the expected fractions $p/(p+q)$ and $q/(p+q)$).

\textbf{Effect of cost-sensitive weighting.}
Under cost-sensitive weighting, same-class neighbors receive higher weight,
so the aggregated representation is dominated by same-class features.
The within-class scatter is bounded by:
\[
  \E[\mathrm{tr}(\mathbf{S}_W)] \leq n\,d\,/\,d_{\mathrm{eff}},
\]
where $d_{\mathrm{eff}} = \E[\sum_j w_{ij}^2]^{-1}$ is an effective degree
that depends on the weight distribution.
Intuitively, concentrating weight on fewer (same-class) neighbors reduces
$d_{\mathrm{eff}}$, which could increase scatter; but this is offset by the
benefit of reducing cross-class contamination in the between-class separation.

\textbf{Why a tight bound is hard.}
The paper notes that a complete formal bound on $\mathrm{tr}(\mathbf{S}_W)$
requires tracking the correlation between the random graph structure and the
cost function (which itself depends on the features).
Since the cost function is learned and feature-dependent, the weights $w_{ij}$
are not independent of the feature noise, making a tight bound technically
challenging.
The key qualitative conclusion is that cost-sensitive weighting improves the
between-class separation (Theorem~2) without catastrophically increasing
within-class scatter, because the noise-averaging effect over $\Theta(n)$
neighbors still applies.


%% ================================================================
\section{Summary of Notation}
\label{sec:notation}
%% ================================================================

For quick reference, here is a table of all major symbols used in the paper:

\begin{center}
\small
\renewcommand{\arraystretch}{1.3}
\begin{tabular}{lll}
\toprule
\textbf{Symbol} & \textbf{Type / Dimensions} & \textbf{Meaning} \\
\midrule
$\G = (V, E)$ & Graph & Undirected graph with nodes $V$ and edges $E$ \\
$n = |V|$ & Scalar (integer) & Number of nodes \\
$d$ & Scalar (integer) & Dimension of node features \\
$d'$ & Scalar (integer) & Dimension of hidden representations \\
$C$ & Scalar (integer) & Number of classes \\
$\bx_i \in \R^d$ & Vector & Raw feature vector of node $i$ \\
$\mathbf{X} \in \R^{n \times d}$ & Matrix & Feature matrix (row $i$ = $\bx_i$) \\
$y_i \in \{1,\ldots,C\}$ & Scalar (integer) & Class label of node $i$ \\
$\mathcal{H}$ & Scalar $\in [0,1]$ & Edge homophily ratio \\
$\N(i)$ & Set of nodes & Neighborhood of $i$ (including $i$) \\
$\bh_i^{(\ell)} \in \R^d$ & Vector & Node $i$'s representation at layer $\ell$ \\
$\bW^{(\ell)} \in \R^{d' \times d}$ & Matrix & Learnable weight matrix at layer $\ell$ \\
$\alpha_{ij}$ & Scalar & Aggregation weight for edge $(i,j)$ \\
$\sigma(\cdot)$ & Function & Nonlinear activation (ReLU or sigmoid) \\
$\bW_g \in \R^{d' \times d}$ & Matrix & Learned projection for cost computation \\
$g_{ij}$ & Scalar $\geq 0$ & Observed divergence (L2 distance in projection) \\
$\tau$ & Scalar $> 0$ & Temperature for sigmoid concordance \\
$s_{ij}$ & Scalar $\in (0,0.5]$ & Concordance score \\
$\tilde{s}_{ij}$ & Scalar $\in (0,1)$ & Softmax-normalized concordance weight \\
$\tilde{d}_{ij}$ & Scalar $\in (0,1)$ & Softmax-normalized discordance weight \\
$\bW_{\text{con}}, \bW_{\text{dis}} \in \R^{d' \times d}$ & Matrices & Channel-specific weight matrices \\
$\bW_{\text{self}} \in \R^{d' \times d}$ & Matrix & Ego (self) channel weight matrix \\
$\gamma_k$ & Scalar $\in [0,1]$ & Gate weight for channel $k$ ($\sum_k \gamma_k = 1$) \\
$\bW_\gamma \in \R^{3 \times 3d'}$ & Matrix & Gate weight matrix \\
$h_{ij}$ & Scalar $\geq 0$ & Learned cost component (extended variant) \\
$\mathbf{a} \in \R^{2d'}$ & Vector & Learnable parameter for $h_{ij}$ \\
$f_{ij}$ & Scalar $\geq 0$ & Total edge cost ($g_{ij} + h_{ij}$, extended variant) \\
$\mathcal{L}_{\text{cal}}$ & Scalar & Calibration regularization loss \\
$\lambda_{\text{cal}}$ & Scalar & Regularization weight (fixed at 0.1) \\
$p$ & Scalar $\in [0,1]$ & CSBM intra-class edge probability \\
$q$ & Scalar $\in [0,1]$ & CSBM inter-class edge probability \\
$\mu$ & Scalar $> 0$ & CSBM feature signal strength \\
$w_+$ & Scalar $\geq 0$ & Expected weight on same-class edges \\
$w_-$ & Scalar $\geq 0$ & Expected weight on different-class edges \\
\bottomrule
\end{tabular}
\end{center}


\end{document}
